\documentclass[11pt,a4paper]{article}
\usepackage[utf8]{inputenc}
\usepackage{amsmath,amssymb,amsthm}
\usepackage{listings}
\usepackage{xcolor}
\usepackage[margin=2.5cm]{geometry}
\usepackage{hyperref}

\lstset{
  language=C++,
  basicstyle=\ttfamily\small,
  keywordstyle=\color{blue}\bfseries,
  commentstyle=\color{green!60!black},
  stringstyle=\color{red!70!black},
  numbers=left,
  numberstyle=\tiny\color{gray},
  breaklines=true,
  frame=single,
  tabsize=4,
  showstringspaces=false
}

\title{IOI 2000 -- Palindrome}
\author{Solution}
\date{}

\begin{document}
\maketitle

\section{Problem Statement}

Given a string $S$ of length $N$ ($N \le 5000$), find the minimum number of
characters that must be inserted (at any positions) to make $S$ a palindrome.

\section{Solution Approach}

\subsection{Key Observation}

The minimum number of insertions equals:
\[
  N - \mathrm{LPS}(S),
\]
where $\mathrm{LPS}(S)$ is the length of the \textbf{longest palindromic subsequence}
of $S$.

\begin{proof}
Every palindrome of length $\ell$ that contains $S$ as a subsequence has length
$\ge 2N - \ell$ where $\ell$ is the longest palindromic subsequence of $S$
(the LPS characters stay in place; each non-LPS character requires one insertion
to create its mirror). The minimum palindrome length achievable is $2N - \mathrm{LPS}(S)$,
requiring $N - \mathrm{LPS}(S)$ insertions.
\end{proof}

\subsection{Reduction to LCS}

The longest palindromic subsequence of $S$ equals the longest common subsequence
of $S$ and its reverse $S^R$:
\[
  \mathrm{LPS}(S) = \mathrm{LCS}(S, S^R).
\]

\subsection{Space-Optimized LCS}

Standard LCS uses an $N \times N$ table, which for $N = 5000$ is 100 MB (too large).
Since each row of the LCS table depends only on the previous row, we use a rolling
two-row approach, reducing space to $O(N)$.

\section{C++ Solution}

\begin{lstlisting}
#include <cstdio>
#include <cstring>
#include <algorithm>
#include <vector>
using namespace std;

int main() {
    int N;
    scanf("%d", &N);

    char S[5001];
    scanf("%s", S);

    // Reverse of S
    char T[5001];
    for (int i = 0; i < N; i++)
        T[i] = S[N - 1 - i];
    T[N] = '\0';

    // LCS(S, T) with O(N) space
    vector<int> prev(N + 1, 0), curr(N + 1, 0);
    for (int i = 1; i <= N; i++) {
        for (int j = 1; j <= N; j++) {
            if (S[i - 1] == T[j - 1]) {
                curr[j] = prev[j - 1] + 1;
            } else {
                curr[j] = max(prev[j], curr[j - 1]);
            }
        }
        swap(prev, curr);
        fill(curr.begin(), curr.end(), 0);
    }

    int lps = prev[N];
    printf("%d\n", N - lps);

    return 0;
}
\end{lstlisting}

\subsection{Alternative: Direct Interval DP}

Define $\mathrm{dp}[i][j]$ = minimum insertions to make $S[i \ldots j]$ a palindrome.

\textbf{Base:} $\mathrm{dp}[i][i] = 0$, $\mathrm{dp}[i][i{-}1] = 0$.

\textbf{Transition:}
\[
  \mathrm{dp}[i][j] = \begin{cases}
    \mathrm{dp}[i{+}1][j{-}1] & \text{if } S[i] = S[j], \\
    \min(\mathrm{dp}[i{+}1][j],\; \mathrm{dp}[i][j{-}1]) + 1 & \text{otherwise}.
  \end{cases}
\]

This also runs in $O(N^2)$ time. With a rolling-array technique (keeping only three
gap-lengths at a time: current, previous, and two-back), the space reduces to $O(N)$.

\begin{lstlisting}
#include <cstdio>
#include <cstring>
#include <algorithm>
#include <vector>
using namespace std;

int main() {
    int N;
    scanf("%d", &N);

    char S[5001];
    scanf("%s", S);

    // Rolling arrays for gap-based interval DP
    // prev2[i] = dp for gap-2 starting at i
    // dp[i]    = dp for gap-1 starting at i
    // ndp[i]   = dp for current gap starting at i
    vector<int> dp(N, 0), ndp(N, 0), prev2(N, 0);

    // gap = 1 (length-2 substrings)
    for (int i = 0; i + 1 < N; i++)
        ndp[i] = (S[i] == S[i + 1]) ? 0 : 1;

    if (N >= 2) {
        swap(prev2, dp); // prev2 = gap-0 results (all 0)
        swap(dp, ndp);   // dp = gap-1 results
    }

    for (int gap = 2; gap < N; gap++) {
        for (int i = 0; i + gap < N; i++) {
            int j = i + gap;
            if (S[i] == S[j])
                ndp[i] = prev2[i + 1]; // dp[i+1][j-1]
            else
                ndp[i] = min(dp[i + 1], dp[i]) + 1;
        }
        swap(prev2, dp);
        swap(dp, ndp);
    }

    printf("%d\n", dp[0]);

    return 0;
}
\end{lstlisting}

\section{Complexity Analysis}

\begin{itemize}
  \item \textbf{Time complexity:} $O(N^2)$ for both approaches.
  \item \textbf{Space complexity:} $O(N)$ with the rolling-array optimization.
\end{itemize}

\end{document}
